\chapter{Topologia dos espaços métricos}\label{ch:b2:metric}

A topologia da reta real (volume do primeiro ano de graduação) generaliza-se, quase
sem mudar uma palavra, a qualquer conjunto munido de uma distância. O
ganho é enorme: sequências de funções, de matrizes, de curvas --- todas
se tornam pontos de \omterm{def:b2:metric:def}{espaços métricos}, e os três pilares aqui demonstrados
--- a \omterm{def:b2:metric:complete}{completude} com o teorema do ponto fixo de Banach, a \omterm{def:b2:metric:compact}{compacidade},
a \omterm{def:b2:metric:connected}{conexidade} --- aplicam-se a elas uniformemente. Este capítulo é a
espinha dorsal de toda a metade analítica do livro.

\section{Espaços métricos}

\begin{definition}\label{def:b2:metric:def}
Um \emph{espaço métrico}\index{espaço métrico} é um conjunto $X$ munido de uma aplicação $d
\colon X \times X \to \R_+$ tal que, para todos $x, y, z$:
\[
d(x,y) = 0 \iff x = y,
\qquad
d(x,y) = d(y,x),
\qquad
d(x,z) \leq d(x,y) + d(y,z).
\]
Bolas: $B(a, r) = \{x : d(a,x) < r\}$ (\omterm{def:b2:metric:topology}{aberta}), $\overline B(a,r) =
\{x : d(a,x) \leq r\}$ (fechada). Uma parte $A \subseteq X$ torna-se um
espaço métrico com a distância induzida.
\end{definition}

\begin{example}\label{ex:b2:metric:examples}
$\R$ com $\abs{x - y}$; $\R^n$ com qualquer uma das distâncias
\[
d_1(x,y) = \sum_i \abs{x_i - y_i},
\quad
d_2(x,y) = \Bigl(\sum_i (x_i - y_i)^2\Bigr)^{1/2},
\quad
d_\infty(x,y) = \max_i \abs{x_i - y_i};
\]
o conjunto $C(\intcc{a}{b})$ das funções \omterm{def:b2:metric:continuity}{contínuas} com a
\emph{distância do sup} $d_\infty(f, g) = \sup_{\intcc{a}{b}} \abs{f -
g}$ (finita: $f - g$ é limitada); qualquer conjunto com a distância
\emph{discreta} ($d(x,y) = 1$ para $x \neq y$). As distâncias provenientes de normas
são o assunto do~\cref{ch:b2:nvs}.
\end{example}

\begin{definition}[Topologia de um espaço métrico]\label{def:b2:metric:topology}
$U \subseteq X$ é \emph{aberto}\index{aberto!num espaço métrico}
quando todo ponto de $U$ é centro de uma bola contida em $U$;
$F$ é \emph{fechado} quando seu complementar é aberto. Vizinhanças,
interior, fecho, densidade e fronteira definem-se exatamente como na
reta real (volume do primeiro ano de graduação), com bolas no lugar de intervalos, e os
enunciados aí demonstrados --- uniões e interseções de abertos,
caracterizações do interior e do fecho, o fecho como menor
fechado que contém o conjunto --- transferem-se com as mesmas demonstrações. As bolas abertas são
abertas, as bolas fechadas são fechadas (desigualdade triangular).
\end{definition}

\begin{example}[Interior, fecho e fronteira num só conjunto]\label{ex:b2:metric:intclosure}
Em $\R$, seja $A = \intoc{0}{1} \cup \{2\}$. Interior:
$\intoo{0}{1}$ --- em torno de qualquer $x \in \intoo01$ uma bola pequena
permanece em $A$; em torno de $1$, toda bola $\intoo{1-r}{1+r}$ vaza
para fora de $A$ pela direita, logo $1$ não é interior; e o
ponto isolado $2$ também não é interior. Fecho:
$\intcc{0}{1} \cup \{2\}$ (o ponto $0$ é limite de $A$,
e nada mais se acrescenta). Fronteira (fecho menos interior):
$\{0, 1, 2\}$. Note as assimetrias que vale a pena guardar: uma
extremidade pode pertencer a um conjunto sem ser interior ($1$), pode
ser aderente sem pertencer ($0$), e um ponto isolado é
sua própria fronteira ($2$). A mesma contabilidade roda palavra por palavra em
qualquer \omterm{def:b2:metric:def}{espaço métrico}, com bolas no lugar de intervalos.
\end{example}

\begin{definition}[Limites, continuidade]\label{def:b2:metric:continuity}
$x_n \to x$ em $X$ quando $d(x_n, x) \to 0$. Uma aplicação $f \colon X \to Y$
entre \omterm{def:b2:metric:def}{espaços métricos} é \emph{contínua}\index{contínua} em $a$ quando
\[
\forall \varepsilon > 0,\ \exists\delta > 0,\quad
d_X(x, a) \leq \delta \implies d_Y\bigl(f(x), f(a)\bigr) \leq
\varepsilon ;
\]
equivalentemente (mesma demonstração que em $\R$), $f(x_n) \to f(a)$ para toda
sequência $x_n \to a$. $f$ é \emph{lipschitziana}\index{lipschitziana} de constante $k$ quando
$d_Y(f(x), f(y)) \leq k\, d_X(x, y)$ sempre --- e então é uniformemente
contínua, logo contínua.
\end{definition}

\begin{theorem}[Caracterização global da continuidade]\label{thm:b2:metric:globalcontinuity}
$f \colon X \to Y$ é \omterm{def:b2:metric:continuity}{contínua} (em todo ponto) se e somente se a
imagem inversa de todo \omterm{def:b2:metric:topology}{aberto} é um \omterm{def:b2:metric:topology}{aberto} --- se e somente se a imagem inversa de todo
fechado é fechada.
\end{theorem}

\begin{proof}
($\Rightarrow$) Seja $V \subseteq Y$ um \omterm{def:b2:metric:topology}{aberto} e $a \in f^{-1}(V)$:
alguma bola $B(f(a), \varepsilon) \subseteq V$; a \omterm{def:b2:metric:continuity}{continuidade} em $a$
fornece $\delta$ com $f(B(a, \delta)) \subseteq B(f(a),
\varepsilon)$, logo $B(a, \delta) \subseteq f^{-1}(V)$.

($\Leftarrow$) Dados $a$ e $\varepsilon$: $f^{-1}\bigl(B(f(a),
\varepsilon)\bigr)$ é \omterm{def:b2:metric:topology}{aberto} e contém $a$, logo contém uma bola
$B(a, \delta)$: essa é a definição de \omterm{def:b2:metric:continuity}{continuidade} em $a$. Conjuntos
fechados: complementares (\cref{prop:b2:structures:images}).
\end{proof}

\section{Espaços completos}

\begin{definition}\label{def:b2:metric:complete}
Uma sequência $(x_n)$ é \emph{de Cauchy} quando $\sup_{p, q \geq N} d(x_p,
x_q) \to 0$ com $N \to \infty$. Um \omterm{def:b2:metric:def}{espaço métrico} é
\emph{completo}\index{espaço métrico completo} quando toda sequência de Cauchy
converge. Convergente $\Rightarrow$ de Cauchy, sempre; as partes fechadas
de espaços completos são completas, e as partes completas de qualquer
espaço são fechadas (mesmas demonstrações que em $\R$: volume do primeiro ano de graduação).
\end{definition}

\begin{example}[De Cauchy sem limite]\label{ex:b2:metric:cauchynolimit}
Em $X = \Q$ com a distância usual, os truncamentos decimais de
$\sqrt2$,
\[
x_0 = 1,\quad x_1 = 1.4,\quad x_2 = 1.41,\quad x_3 = 1.414,
\quad\dots
\]
satisfazem $\abs{x_p - x_q} \leq 10^{-\min(p,q)}$: são de Cauchy em $\Q$.
Um limite em $\Q$ seria também o limite em $\R$, a saber
$\sqrt2 \notin \Q$: não existe limite em $X$. A incompletude é
a presença desses ``limites fantasmas''; a \omterm{def:b2:metric:complete}{completude} de $\R$
foi projetada no volume do primeiro ano de graduação precisamente para dar a toda
sequência de Cauchy um lar.
\end{example}

\begin{theorem}\label{thm:b2:metric:rncomplete}
$\R^n$ (com qualquer uma das três distâncias do
\cref{ex:b2:metric:examples}) e $\bigl(C(\intcc{a}{b}),
d_\infty\bigr)$ são \omterm{def:b2:metric:complete}{completos}.
\end{theorem}

\begin{proof}
$\R^n$: uma sequência de Cauchy é de Cauchy em cada coordenada (cada
$\abs{x_i - y_i} \leq d(x,y)$ para as três distâncias), logo cada
coordenada converge (\omterm{def:b2:metric:complete}{completude} de $\R$, volume do primeiro ano de graduação), e a
convergência coordenada a coordenada implica a convergência para $d_\infty$
(um número finito de coordenadas), logo para as três (as três
distâncias se dominam mutuamente a menos de fatores constantes: $d_\infty
\leq d_2 \leq d_1 \leq n\,d_\infty$).

$C(\intcc{a}{b})$: seja $(f_n)$ de Cauchy para $d_\infty$. Para cada $x$,
$(f_n(x))$ é de Cauchy em $\R$ ($\abs{f_p(x) - f_q(x)} \leq
d_\infty(f_p, f_q)$): converge para algum $f(x)$. Passando ao
limite em $\abs{f_p(x) - f_q(x)} \leq \varepsilon$ (válido para $p, q
\geq N_\varepsilon$, todos os $x$) quando $q \to \infty$: $\abs{f_p(x) -
f(x)} \leq \varepsilon$ para todo $x$, isto é, $d_\infty(f_p, f) \leq
\varepsilon$: convergência uniforme. O limite é \omterm{def:b2:metric:continuity}{contínuo}: dado
$\varepsilon$, escolha $p$ com $\sup\abs{f_p - f} \leq \varepsilon$,
e depois use a \omterm{def:b2:metric:continuity}{continuidade} de $f_p$ em $a$ e a decomposição em três termos
\[
\abs{f(x) - f(a)} \leq \abs{f(x) - f_p(x)} + \abs{f_p(x) - f_p(a)}
+ \abs{f_p(a) - f(a)} \leq 3\varepsilon
\]
para $x$ próximo de $a$. (Esse ``argumento dos $3\varepsilon$'' reaparece como
o teorema do limite uniforme do~\cref{ch:b2:funcseq}.)
\end{proof}

\begin{example}[Abertos e fechados reconhecidos pela continuidade]\label{ex:b2:metric:recognize}
A caracterização global
(\cref{thm:b2:metric:globalcontinuity}) é a ferramenta cotidiana da
contabilidade topológica. Em $\R^2$: o conjunto $\{(x, y) : x^2 + y^2
< 1,\ y > x^3\}$ é \omterm{def:b2:metric:topology}{aberto} --- é $g^{-1}(\intoo{-\infty}{1})
\cap h^{-1}(\intoo{0}{+\infty})$ para a função \omterm{def:b2:metric:continuity}{contínua} $g(x,y) =
x^2 + y^2$ e $h(x, y) = y - x^3$, uma interseção de duas imagens inversas de
\omterm{def:b2:metric:topology}{abertos}. Em $\bigl(C(\intcc01), d_\infty\bigr)$: o conjunto das
funções com $f(0) = f(1)$ e $\int_0^1 f = 0$ é fechado ---
é a imagem inversa de $\{(0,0)\}$ pela aplicação \omterm{def:b2:metric:continuity}{contínua} $f \mapsto
\bigl(f(0) - f(1),\ \int_0^1 f\bigr)$ com valores em $\R^2$ (cada
coordenada é $1$-lipschitziana, como no~\cref{exo:b2:metric:3}). O
método nunca desenha figura: exiba uma aplicação \omterm{def:b2:metric:continuity}{contínua}, leia o
conjunto como imagem inversa, cite o teorema.
\end{example}

\begin{example}[Um fechado definido por uma infinidade de condições]\label{ex:b2:metric:lipclosed}
Em $\bigl(C(\intcc01), d_\infty\bigr)$, o conjunto
\[
L = \{f : \abs{f(x) - f(y)} \leq \abs{x - y}
\ \text{para todos } x, y\}
\]
das funções $1$-lipschitzianas é fechado, embora seja recortado por
uma quantidade não \omterm{def:b2:structures:countable}{enumerável} de condições: para cada par fixo $(x, y)$, a
aplicação $f \mapsto \abs{f(x) - f(y)} - \abs{x - y}$ é \omterm{def:b2:metric:continuity}{contínua}
(as avaliações são $1$-lipschitzianas), de modo que cada condição isolada
define um fechado, e $L$ é a \emph{interseção} dessa
família --- e uma interseção arbitrária de fechados é
fechada. O mesmo molde certifica que são fechados os conjuntos de funções
monótonas, de funções convexas, de funções limitadas por um $g$ fixo:
os limites uniformes herdam toda propriedade que se exprima como família
de restrições pontuais fechadas. O que os limites uniformes \emph{não}
herdam automaticamente --- a derivabilidade, por exemplo --- é exatamente
aquilo pelo que o~\cref{ch:b2:funcseq} terá de penar.
\end{example}

\begin{theorem}[Teorema do ponto fixo de Banach]\label{thm:b2:metric:banach}
Seja $X$ um \omterm{def:b2:metric:complete}{espaço métrico completo} não vazio e $f \colon X \to X$ uma
\emph{contração}: \omterm{def:b2:metric:continuity}{lipschitziana} de constante $k < 1$. Então $f$ tem um
único ponto fixo $\ell$, e toda órbita $x_{n+1} = f(x_n)$
converge para $\ell$, com\index{teorema do ponto fixo de Banach}
\[
d(x_n, \ell) \leq \frac{k^n}{1 - k}\, d(x_1, x_0) .
\]
\end{theorem}

\begin{proof}
Unicidade: dois pontos fixos estão a uma distância $\leq k$ vezes ela mesma.
Existência: $d(x_{n+1}, x_n) \leq k^n d(x_1, x_0)$ por indução, logo,
para $q > p$,
\[
d(x_q, x_p) \leq \sum_{j=p}^{q-1} d(x_{j+1}, x_j)
\leq d(x_1, x_0) \sum_{j \geq p} k^j
= \frac{k^p}{1-k}\, d(x_1, x_0) \xrightarrow[p\to\infty]{} 0 :
\]
de Cauchy, logo convergente para algum $\ell$; a \omterm{def:b2:metric:continuity}{continuidade} de $f$ passa
$x_{n+1} = f(x_n)$ ao limite: $\ell = f(\ell)$. A cota de erro
é a estimativa exibida com $q \to \infty$.
\end{proof}

\begin{example}[Uma equação integral]\label{ex:b2:metric:integralequation}
Em $X = C(\intcc{0}{1})$ (\omterm{def:b2:metric:complete}{completo},
\cref{thm:b2:metric:rncomplete}), considere $T(f)(x) = 1 + \frac12
\int_0^x f(t)\,\dd t$. Para $f, g \in X$:
\[
\abs{T(f)(x) - T(g)(x)} \leq \frac12 \int_0^x \abs{f - g} \leq
\frac12\, d_\infty(f, g),
\]
de modo que $T$ é uma $\frac12$-contração: ela tem um único ponto fixo
\omterm{def:b2:metric:continuity}{contínuo} --- a solução de $f' = \frac f2$, $f(0) = 1$, a saber
$\eu^{x/2}$. Esse esquema, industrializado, torna-se o
teorema de Cauchy--Lipschitz do~\cref{ch:b2:diffeq}.
\end{example}

\begin{example}[Um ponto fixo numérico: $x = \cos x$]\label{ex:b2:metric:cosfix}
No \omterm{def:b2:metric:complete}{completo} $X = \intcc{0}{1}$, a aplicação $f = \cos$ leva $X$
em $\intcc{\cos 1}{1} \subseteq X$ e é uma contração: pela
desigualdade do valor médio,
\[
\abs{\cos x - \cos y} \leq \bigl(\sup_{\intcc01}\abs{\sin}\bigr)
\abs{x - y} = (\sin 1)\abs{x - y},
\qquad \sin 1 \approx 0.841 < 1 .
\]
Banach: uma única solução de $x = \cos x$ em $\intcc01$ (logo em
$\R$: todo ponto fixo real está em $\intcc{-1}{1}$ e depois em
$\intcc{\cos 1}{1}$ após uma aplicação), e a iteração
$x_{n+1} = \cos x_n$ converge para ela a partir de qualquer início: $x_\infty
\approx 0.739085$, o famoso número obtido martelando a
tecla do cosseno de uma calculadora. A cota de erro prevê
um decaimento em $(\sin1)^n/(1 - \sin1)$ --- cerca de um algarismo a cada $13$
toques; a cota a posteriori do problema de fim de semana deste
capítulo (questão 14) certifica cada passo em tempo real.
\end{example}

\section{Compacidade}

\begin{definition}\label{def:b2:metric:compact}
Um \omterm{def:b2:metric:def}{espaço métrico} $X$ é \emph{compacto}\index{espaço métrico compacto}
quando toda sequência em $X$ tem uma subsequência que converge \emph{em}
$X$ (a propriedade de Bolzano--Weierstrass). Uma parte é compacta quando o
é com a distância induzida.
\end{definition}

\begin{theorem}[Primeiras propriedades]\label{thm:b2:metric:compactprops}
\begin{enumerate}
  \item Uma parte \omterm{def:b2:metric:compact}{compacta} é fechada e limitada; uma parte fechada de um
        espaço \omterm{def:b2:metric:compact}{compacto} é \omterm{def:b2:metric:compact}{compacta}.
  \item Em $\R^n$, vale a recíproca: \omterm{def:b2:metric:compact}{compacta} $\iff$ fechada e
        limitada.
  \item A imagem \omterm{def:b2:metric:continuity}{contínua} de um espaço \omterm{def:b2:metric:compact}{compacto} é \omterm{def:b2:metric:compact}{compacta}; uma função
        real \omterm{def:b2:metric:continuity}{contínua} num espaço \omterm{def:b2:metric:compact}{compacto} não vazio é
        limitada e atinge seus extremos.
  \item (Heine) Uma aplicação \omterm{def:b2:metric:continuity}{contínua} num espaço \omterm{def:b2:metric:compact}{compacto} é uniformemente
        contínua.
  \item Produtos: se $X, Y$ são \omterm{def:b2:metric:compact}{compactos}, então $X \times Y$ também é (com
        $d\bigl((x,y),(x',y')\bigr) = d(x,x') + d(y,y')$).
\end{enumerate}
\end{theorem}

\begin{proof}
(1) Mesmos argumentos que na reta (volume do primeiro ano de graduação): uma sequência que escapa para o%
infinito, ou que converge para fora, não tem subsequência
que convirja dentro; para a segunda afirmação, extraia no
\omterm{def:b2:metric:compact}{compacto} ambiente e use que o conjunto é fechado.

(2) As sequências limitadas de $\R^n$ têm subsequências convergentes
componente a componente: extraia na primeira coordenada
(Bolzano--Weierstrass em $\R$), depois, dessa subsequência, na
segunda, e assim por diante ($n$ extrações sucessivas); ser fechado mantém
o limite dentro.

(3) Dada $(f(x_n))$, extraia $x_{\varphi(n)} \to x \in X$;
a \omterm{def:b2:metric:continuity}{continuidade} dá $f(x_{\varphi(n)}) \to f(x) \in f(X)$. Caso real: a
\omterm{def:b2:metric:compact}{compacidade} de $f(X) \subseteq \R$ o torna fechado e limitado, e
$\sup f(X) \in f(X)$ (o supremo de um conjunto é aderente a ele, e
$f(X)$ é fechado).

(4) A demonstração do primeiro ano se transfere palavra por palavra; ei-la, em roupagem
métrica. Suponha $f \colon X \to Y$ \omterm{def:b2:metric:continuity}{contínua} no \omterm{def:b2:metric:compact}{compacto} $X$
mas não uniformemente contínua: algum $\varepsilon > 0$ admite, para
cada $n$, pontos com
\[
d_X(x_n, y_n) \leq \frac{1}{n+1}
\qquad\text{e}\qquad
d_Y\bigl(f(x_n), f(y_n)\bigr) > \varepsilon .
\]
Extraia $x_{\varphi(n)} \to a \in X$; então $y_{\varphi(n)} \to
a$ também (as distâncias mútuas tendem a $0$). A \omterm{def:b2:metric:continuity}{continuidade} em $a$
manda as duas sequências de imagens para $f(a)$, logo
$d_Y\bigl(f(x_{\varphi(n)}), f(y_{\varphi(n)})\bigr) \to 0$ ---
contradizendo o afastamento uniforme $> \varepsilon$. A \omterm{def:b2:metric:compact}{compacidade}
forneceu exatamente uma coisa: o ponto de acumulação $a$ no qual
aplicar a \omterm{def:b2:metric:continuity}{continuidade} simples.

(5) Extraia nas coordenadas de $X$, depois de novo nas
coordenadas de $Y$.
\end{proof}

\begin{example}[O teorema de Heine, com e sem compacidade]\label{ex:b2:metric:heinework}
Em $\intcc{0}{1}$, a função $x \mapsto x^2$ é uniformemente
contínua --- Heine o afirma sem cálculo algum, mas a estimativa
direta é instrutiva:
\[
\abs{x^2 - y^2} = \abs{x + y}\,\abs{x - y} \leq 2\abs{x - y},
\]
de modo que $\delta = \varepsilon/2$ serve \emph{para todos os pontos de uma
vez}. Em $\R$ a mesma função não é uniformemente contínua:
com $x_n = n$ e $y_n = n + \frac1n$, o afastamento $\abs{x_n - y_n}
= \frac1n \to 0$ enquanto $\abs{x_n^2 - y_n^2} = 2 + \frac1{n^2}
\geq 2$: nenhum $\delta$ único serve para $\varepsilon = 1$. O
mecanismo é visível: a constante \omterm{def:b2:metric:continuity}{lipschitziana} local $\abs{x + y}$
é limitada num \omterm{def:b2:metric:compact}{compacto} e ilimitada em $\R$ --- o teorema de
Heine é exatamente a afirmação de que a \omterm{def:b2:metric:compact}{compacidade} limita tais
constantes locais uniformemente.
\end{example}

\begin{method}[Demonstrar que um conjunto é compacto]\label{met:b2:metric:compactness}
Três caminhos, em ordem de frequência. (1) \emph{Reconhecimento
no ambiente:} em $\R^n$ (ou em qualquer espaço normado de
dimensão finita, \cref{ch:b2:nvs}), verifique que é fechado --- tipicamente como
imagem inversa, \cref{ex:b2:metric:recognize} --- e limitado. (2)
\emph{Herança:} uma parte fechada de um \omterm{def:b2:metric:compact}{compacto} conhecido é
\omterm{def:b2:metric:compact}{compacta}; uma união finita ou um produto de \omterm{def:b2:metric:compact}{compactos} é \omterm{def:b2:metric:compact}{compacto}; a
imagem \omterm{def:b2:metric:continuity}{contínua} de um \omterm{def:b2:metric:compact}{compacto} é \omterm{def:b2:metric:compact}{compacta}. (3) \emph{Na
unha:} extraia uma subsequência convergente de uma sequência
arbitrária --- em geral por extrações sucessivas coordenada a
coordenada. Para provar a \emph{não} \omterm{def:b2:metric:compact}{compacidade}, basta uma testemunha:
uma sequência sem subsequência convergente, quase sempre
pontos a distância mútua $\geq \varepsilon$.
\end{method}

\begin{example}[Distâncias entre conjuntos: a compacidade ganha o pão]\label{ex:b2:metric:setdistance}
Sejam $K$ \omterm{def:b2:metric:compact}{compacto}, $F$ fechado, $K \cap F = \emptyset$ num
\omterm{def:b2:metric:def}{espaço métrico}. Então
\[
d(K, F) = \inf\,\{d(x, y) : x \in K,\ y \in F\} > 0 :
\]
a função $x \mapsto d(x, F)$ é \omterm{def:b2:metric:continuity}{contínua}
(\cref{exo:b2:metric:11}) e positiva em $K$ ($d(x, F) = 0$
poria $x \in \overline F = F$), logo atinge um mínimo positivo
no \omterm{def:b2:metric:compact}{compacto} $K$
(\cref{thm:b2:metric:compactprops}~(3)). A \omterm{def:b2:metric:compact}{compacidade} não é
decorativa: para dois \emph{fechados} o ínfimo pode se anular
sem ser atingido --- em $\R^2$, a hipérbole $F_1 = \{xy
= 1\}$ e o eixo $F_2 = \{y = 0\}$ são fechados disjuntos
com $d(F_1, F_2) = 0$ (os pontos $(n, \frac1n)$ se aproximam do
eixo). A fuga para o infinito é exatamente o que a \omterm{def:b2:metric:compact}{compacidade} proíbe.
\end{example}

\begin{theorem}[Borel--Lebesgue]\label{thm:b2:metric:borellebesgue}
Um \omterm{def:b2:metric:def}{espaço métrico} $X$ é \omterm{def:b2:metric:compact}{compacto} se e somente se toda cobertura de $X$ por
\omterm{def:b2:metric:topology}{abertos} admite uma subcobertura \emph{finita}.\index{teorema de
Borel--Lebesgue}
\end{theorem}

\begin{proof}
($\Leftarrow$) Suponha que $(x_n)$ não tenha subsequência convergente. Afirmamos
que todo $x \in X$ tem uma bola $B(x, r_x)$ que contém $x_n$
apenas para um número finito de índices $n$: caso contrário, toda bola
$B(x, \frac1{k+1})$ conteria uma infinidade de termos e,
escolhendo índices
\[
\varphi(0) < \varphi(1) < \varphi(2) < \cdots
\quad\text{com}\quad
x_{\varphi(k)} \in B\Bigl(x, \frac{1}{k+1}\Bigr)
\]
(possível em cada passo precisamente porque restam infinitos
candidatos), construiríamos uma subsequência convergindo para $x$.
As bolas $B(x, r_x)$ cobrem $X$; se um número finito delas
cobrisse $X$, o conjunto de índices $\N$ seria uma união finita de
conjuntos finitos: absurdo.

($\Rightarrow$) Duas etapas. \emph{Número de Lebesgue:} para uma cobertura \omterm{def:b2:metric:topology}{aberta}
$(U_i)$ de um \omterm{def:b2:metric:compact}{compacto} $X$, existe $\rho > 0$ tal que toda bola
de raio $\rho$ está contida em algum $U_i$. Caso contrário, para cada $n$ escolha
$x_n$ com $B(x_n, \frac{1}{n+1})$ contida em nenhum $U_i$; extraia
$x_{\varphi(n)} \to x \in U_{i_0} \supseteq B(x, r)$; para $n$ grande,
$B(x_{\varphi(n)}, \frac{1}{\varphi(n)+1}) \subseteq B(x, r)
\subseteq U_{i_0}$: contradição. \emph{Precompacidade:} para
todo $\rho > 0$, um número finito de bolas de raio $\rho$ cobre $X$.
Caso contrário, escolha indutivamente $x_{n+1}$ fora de $B(x_0, \rho) \cup
\dots \cup B(x_n, \rho)$: a sequência tem distâncias mútuas $\geq
\rho$, logo nenhuma subsequência de Cauchy --- e portanto nenhuma convergente:
contradição. Combinando: cubra $X$ por um número finito de bolas de raio
$\rho$ (o número de Lebesgue), cada uma dentro de algum $U_i$: uma
subcobertura finita.
\end{proof}

\begin{example}[Uma $\varepsilon$-rede, contada]\label{ex:b2:metric:epsnet}
A precompacidade (extraída da demonstração do
\cref{thm:b2:metric:borellebesgue}) é bem concreta em
$\intcc{0}{1}$: para $\varepsilon > 0$, as $\lceil
\frac{1}{2\varepsilon}\rceil$ bolas centradas em $\varepsilon,
3\varepsilon, 5\varepsilon, \dots$ de raio $\varepsilon$
o cobrem --- cerca de $\frac1{2\varepsilon}$ bolas, e nenhuma cobertura
consegue usar menos de $\frac{1}{2\varepsilon}$ delas (cada
bola cobre comprimento no máximo $2\varepsilon$). Em $\intcc01^2$
a contagem eleva-se ao quadrado, da ordem de $\varepsilon^{-2}$: os números de recobrimento
crescem como $\varepsilon^{-d}$ em dimensão $d$ --- uma face
quantitativa da \omterm{def:b2:metric:compact}{compacidade}, e a razão pela qual as bolas unitárias
de dimensão infinita de \cref{ch:b2:nvs} (onde não existe nenhuma
$\frac13$-rede finita) não podem ser \omterm{def:b2:metric:compact}{compactas}.
\end{example}

\begin{example}[Lendo a compacidade nas coberturas]\label{ex:b2:metric:covers}
O intervalo \omterm{def:b2:metric:topology}{semiaberto} $\intoc{0}{1}$ é coberto pelos \omterm{def:b2:metric:topology}{abertos}
$U_n = \intoo{\frac1n}{2}$, $n \geq 1$; qualquer subfamília finita tem um
maior índice $N$ e deixa de fora $\intoc{0}{\frac1N}$: não há subcobertura
finita, logo $\intoc{0}{1}$ não é \omterm{def:b2:metric:compact}{compacto} --- o que a
definição sequencial enxerga por meio de $x_n = \frac1n$, cujo limite
$0$ escapa. Por outro lado, acrescentar o único ponto $0$
repara os dois diagnósticos de uma vez: em $\intcc{0}{1}$ toda cobertura desse tipo
deve conter um conjunto que contém $0$, o qual engole todo um segmento
inicial, e um número finito de conjuntos termina o resto. As duas linguagens
do~\cref{thm:b2:metric:borellebesgue} sempre falham ou funcionam
juntas --- as coberturas detectam a fuga exatamente onde as sequências detectam.
\end{example}

\begin{remark}[Perspectivas dentro deste volume]\label{rem:b2:metric:perspectives}
Este capítulo é a parede de sustentação do volume; observe onde
cada pilar carrega peso. \emph{\omterm{def:b2:metric:complete}{Completude}}: o critério de
Cauchy torna-se o teste de convergência para séries em espaços de
Banach (\cref{ch:b2:series}), a convergência uniforme do
\cref{ch:b2:funcseq} é exatamente a convergência no \omterm{def:b2:metric:complete}{completo}
$\bigl(C, d_\infty\bigr)$, e Cauchy--Lipschitz
(\cref{ch:b2:diffeq}) é o teorema do ponto fixo de Banach vestido
de equação integral. \emph{\omterm{def:b2:metric:compact}{Compacidade}}: ela demonstra a
equivalência das normas (\cref{ch:b2:nvs}), a obtenção dos
extremos para a otimização do~\cref{ch:b2:diffcalc} e a
existência das melhores aproximações (problema de fim de semana do
\cref{ch:b2:nvs}). \emph{\omterm{def:b2:metric:connected}{Conexidade}}: ela globaliza enunciados locais
--- a unicidade das soluções de equações diferenciais, o
teorema do valor intermediário sobre curvas (\cref{ch:b2:curves}) e
as duas componentes de $GL_n(\R)$ que a teoria da orientação
(\cref{ch:b2:multint}) manterá separadas.
\end{remark}

\begin{remark}[Armadilhas comuns]\label{rem:b2:metric:pitfalls}
(i) ``Fechado e limitado implica \omterm{def:b2:metric:compact}{compacto}'' é um teorema sobre
$\R^n$, não sobre \omterm{def:b2:metric:def}{espaços métricos}: um conjunto infinito com a
métrica discreta é fechado e limitado em si mesmo e no entanto não é \omterm{def:b2:metric:compact}{compacto}
(\cref{exo:b2:metric:4}), e a bola unitária fechada de
$C(\intcc01)$ também falha (\cref{ch:b2:nvs}). (ii) A \omterm{def:b2:metric:complete}{completude} é
uma propriedade da \emph{distância}, não da topologia: $\R$
com $d(x,y) = \abs{\arctan x - \arctan y}$ tem as mesmas sequências
convergentes de sempre, mas é incompleto
(\cref{exo:b2:metric:1}). (iii) Uma bijeção \omterm{def:b2:metric:continuity}{contínua} não precisa
ser um homeomorfismo --- a parametrização do círculo do
\cref{exo:b2:metric:7}; a \omterm{def:b2:metric:compact}{compacidade} da fonte conserta isso.
(iv) O teorema de Banach precisa de $k < 1$ \emph{uniformemente}: a
condição $d(f(x), f(y)) < d(x,y)$ sozinha nada garante num espaço
não \omterm{def:b2:metric:compact}{compacto} (\cref{exo:b2:metric:5}). (v) \omterm{def:b2:metric:connected}{Conexo} não
implica \omterm{def:b2:metric:connected}{conexo} por caminhos em geral --- mas para os \omterm{def:b2:metric:topology}{abertos}
de espaços normados encontrados neste livro, as duas noções coincidem
(\cref{ch:b2:nvs}).
\end{remark}

\begin{remark}[Onde este capítulo é usado]
Em toda parte na metade analítica. A \omterm{def:b2:metric:complete}{completude} de
$C(\intcc{a}{b})$ move os teoremas de convergência do
\cref{ch:b2:funcseq} e a teoria de Cauchy--Lipschitz do
\cref{ch:b2:diffeq} (o problema de fim de semana deste capítulo já demonstra o
teorema local de Picard--Lindelöf); a \omterm{def:b2:metric:compact}{compacidade} dá a
equivalência das normas em dimensão finita (\cref{ch:b2:nvs}) e
a existência de extremos no~\cref{ch:b2:diffcalc}; a \omterm{def:b2:metric:connected}{conexidade}
sustenta os argumentos de valor intermediário do~\cref{ch:b2:realfun}
e a globalização da unicidade para equações diferenciais. No
volume do terceiro ano de graduação, a \omterm{def:b2:metric:compact}{compacidade} em espaços de funções (o
teorema de Arzelà--Ascoli) e o teorema de Baire
(\cref{exo:b2:metric:12} aqui) tornam-se ferramentas de uso diário.
\end{remark}

\begin{omfigure}
\begin{tikzpicture}[xscale=6.4, yscale=0.8]
  \draw[omDef, line width=3pt] (0,3) -- (1,3);
  \draw[omDef, line width=3pt] (0,2) -- (0.3333,2);
  \draw[omDef, line width=3pt] (0.6667,2) -- (1,2);
  \foreach \a in {0, 0.6667} {
    \draw[omDef, line width=3pt] (\a,1) -- (\a + 0.1111,1);
    \draw[omDef, line width=3pt]
      (\a + 0.2222,1) -- (\a + 0.3333,1);
  }
  \foreach \a in {0, 0.2222, 0.6667, 0.8889} {
    \draw[omThm, line width=3pt] (\a,0) -- (\a + 0.037,0);
    \draw[omThm, line width=3pt]
      (\a + 0.0741,0) -- (\a + 0.1111,0);
  }
  \node[left, font=\small] at (0,3) {$C_0$};
  \node[left, font=\small] at (0,2) {$C_1$};
  \node[left, font=\small] at (0,1) {$C_2$};
  \node[left, font=\small] at (0,0) {$C_3$};
\end{tikzpicture}

{\small As primeiras etapas do conjunto de Cantor
(\cref{exo:b2:metric:8}): cada nível apaga o terço médio \omterm{def:b2:metric:topology}{aberto}
de cada segmento. A interseção $C = \bigcap_n C_n$ é
\omterm{def:b2:metric:compact}{compacta}, tem interior vazio e comprimento zero, e no entanto é \omterm{def:b2:structures:countable}{equipotente}
a $\R$ --- e ela reaparece como \emph{ponto fixo} de uma
contração sobre conjuntos no problema de fim de semana deste capítulo (questão
22).}
\end{omfigure}

\section{Conexidade}

\begin{definition}\label{def:b2:metric:connected}
$X$ é \emph{conexo}\index{espaço conexo} quando não admite
partição em dois \omterm{def:b2:metric:topology}{abertos} não vazios --- equivalentemente, quando suas
únicas partes simultaneamente \omterm{def:b2:metric:topology}{abertas} e fechadas são $\emptyset$ e $X$. $X$ é
\emph{conexo por caminhos} quando quaisquer dois pontos são ligados por uma aplicação
\omterm{def:b2:metric:continuity}{contínua} $\gamma \colon \intcc{0}{1} \to X$.
\end{definition}

\begin{theorem}\label{thm:b2:metric:connectedness}
\begin{enumerate}
  \item As partes \omterm{def:b2:metric:connected}{conexas} de $\R$ são exatamente os intervalos.
  \item A imagem \omterm{def:b2:metric:continuity}{contínua} de um \omterm{def:b2:metric:connected}{espaço conexo} é \omterm{def:b2:metric:connected}{conexa} ---
        donde o teorema do valor intermediário geral: uma função
        real \omterm{def:b2:metric:continuity}{contínua} num \omterm{def:b2:metric:connected}{espaço conexo} assume todo valor entre
        dois quaisquer de seus valores.
  \item \omterm{def:b2:metric:connected}{Conexo} por caminhos $\Rightarrow$ \omterm{def:b2:metric:connected}{conexo}. (A recíproca é falsa
        em geral; ela vale para os \omterm{def:b2:metric:topology}{abertos} de espaços normados,
        \cref{ch:b2:nvs}.)
\end{enumerate}
\end{theorem}

\begin{proof}
(1) Um $A$ que não é intervalo deixa de fora algum $z$ situado entre dois de seus pontos:
$A = (A \cap \intoo{-\infty}{z}) \cup (A \cap \intoo{z}{+\infty})$
o reparte em duas peças não vazias e \omterm{def:b2:metric:topology}{abertas} (em $A$). Reciprocamente, seja
$I$ um intervalo e $I = U \cup V$ uma partição em conjuntos não vazios
relativamente \omterm{def:b2:metric:topology}{abertos}; escolha $a \in U$, $b \in V$, digamos $a < b$, e
ponha $s = \sup\,(U \cap \intcc{a}{b})$, um ponto de $\intcc{a}{b}
\subseteq I$. Se $s \in U$: então $s \neq b$, e a abertura relativa
de $U$ põe todo um intervalo em torno de $s$ (intersectado com $I$)
dentro de $U$ --- de modo que há pontos de $U \cap \intcc{a}{b}$ maiores que $s$,
contradizendo o supremo. Se $s \in V$: a abertura relativa de $V$
põe um intervalo $\intoo{s - r}{s + r} \cap I$ dentro de $V$; mas o
supremo é aderente a $U \cap \intcc{a}{b}$, que deve encontrar esse
intervalo --- contradição com $U \cap V = \emptyset$. (Esse é
o argumento aberto-fechado do primeiro ano para $\R$, rodado dentro de $I$.)

(2) Se $f(X) = U' \cup V'$ se reparte em conjuntos não vazios relativamente
\omterm{def:b2:metric:topology}{abertos}, então $X = f^{-1}(U') \cup f^{-1}(V')$ reparte $X$
(\cref{thm:b2:metric:globalcontinuity}). TVI: $f(X) \subseteq \R$
é \omterm{def:b2:metric:connected}{conexo}, logo um intervalo por (1).

(3) Suponha $X = U \cup V$, ambos \omterm{def:b2:metric:topology}{abertos} não vazios, e ligue $a \in U$
a $b \in V$ por um caminho $\gamma$: então $\gamma^{-1}(U),
\gamma^{-1}(V)$ repartem $\intcc{0}{1}$, contradizendo (1).
\end{proof}

\begin{example}\label{ex:b2:metric:glnr}
$GL_n(\R)$ não é \omterm{def:b2:metric:connected}{conexo}: $\det$ é \omterm{def:b2:metric:continuity}{contínuo} (um polinômio nas
entradas) sobre $\R^*$, que não é \omterm{def:b2:metric:connected}{conexo}; as imagens inversas de
$\R_+^*$ e de $\R_-^*$ repartem $GL_n(\R)$. (Cada peça é de fato
\omterm{def:b2:metric:connected}{conexa} por caminhos --- um exercício agradável, além das nossas necessidades.) Em
contraste, $GL_n(\C)$ \emph{é} \omterm{def:b2:metric:connected}{conexo} por caminhos:
\cref{exo:b2:metric:10}.
\end{example}

\begin{example}[Um ponto fixo só pela conexidade]\label{ex:b2:metric:ivtfixed}
Toda $f \colon \intcc01 \to \intcc01$ \omterm{def:b2:metric:continuity}{contínua} tem um ponto
fixo --- sem hipótese de contração, sem iteração. Considere
$g(x) = f(x) - x$, \omterm{def:b2:metric:continuity}{contínua} no \omterm{def:b2:metric:connected}{conexo} $\intcc01$:
\[
g(0) = f(0) \geq 0,
\qquad
g(1) = f(1) - 1 \leq 0 ,
\]
e o teorema do valor intermediário
(\cref{thm:b2:metric:connectedness}~(2)) entrega um zero de
$g$, isto é, um ponto fixo de $f$. Contraste com Banach
(\cref{thm:b2:metric:banach}): aqui a existência é topológica
e gratuita, mas a unicidade e o algoritmo se perdem --- $f =
\mathrm{id}$ tem todo ponto fixo, e a iteração de uma
$f$ não contrativa pode ciclar para sempre. Os dois teoremas
de ponto fixo deste capítulo respondem a perguntas diferentes com
moedas diferentes.
\end{example}

\begin{example}[$\R$ e $\R^2$ não são homeomorfos]\label{ex:b2:metric:rnotr2}
A \omterm{def:b2:metric:connected}{conexidade} é uma impressão digital topológica. Suponha que $h \colon
\R^2 \to \R$ fosse um homeomorfismo (uma bijeção \omterm{def:b2:metric:continuity}{contínua} de inversa
\omterm{def:b2:metric:continuity}{contínua}). Retire um ponto $a \in \R^2$: a
restrição $h \colon \R^2\setminus\{a\} \to
\R\setminus\{h(a)\}$ continua sendo um homeomorfismo. Mas $\R^2$ menos
um ponto é \omterm{def:b2:metric:connected}{conexo} por caminhos --- ligue dois pontos quaisquer por um segmento,
desviando por um segundo segmento através de um ponto auxiliar se
$a$ bloquear o caminho direto --- logo \omterm{def:b2:metric:connected}{conexo}
(\cref{thm:b2:metric:connectedness}~(3)); ao passo que $\R$ menos um
ponto se reparte em duas semirretas \omterm{def:b2:metric:topology}{abertas} não vazias: não é \omterm{def:b2:metric:connected}{conexo}.
A \omterm{def:b2:metric:connected}{conexidade} é preservada por aplicações \omterm{def:b2:metric:continuity}{contínuas}: contradição.
O plano e a reta são genuinamente diferentes \emph{como
espaços topológicos} --- um fato que a cardinalidade sozinha
(bijeções ao estilo \cref{exo:b2:structures:3} realmente existem!) é grosseira
demais para enxergar.
\end{example}

\section{Exercícios}

\begin{exercise}[$\star$]\label{exo:b2:metric:1}
Em $\R$, verifique que $\delta(x, y) = \min(1, \abs{x - y})$ e
$d(x,y) = \abs{\arctan x - \arctan y}$ são distâncias. Quais
sequências convergem para cada uma? $(\R, d)$ é \omterm{def:b2:metric:complete}{completo}?
\end{exercise}

\begin{exercise}[$\star$]\label{exo:b2:metric:2}
Num \omterm{def:b2:metric:def}{espaço métrico}, prove que uma sequência convergente é de Cauchy e
limitada, e que uma sequência de Cauchy com uma subsequência convergente
converge. Deduza de novo que os \omterm{def:b2:metric:compact}{espaços métricos compactos} são \omterm{def:b2:metric:complete}{completos}.
\end{exercise}

\begin{exercise}[$\star$]\label{exo:b2:metric:3}
Em $\bigl(C(\intcc{0}{1}), d_\infty\bigr)$, calcule a distância
entre $f(x) = x$ e $g(x) = x^2$; descreva a bola fechada
$\overline B(0, 1)$; e prove que o conjunto $\{f : f(0) = 0\}$ é
fechado, enquanto $\{f : f(0) > 0\}$ é \omterm{def:b2:metric:topology}{aberto}.
\end{exercise}

\begin{exercise}[$\star\star$]\label{exo:b2:metric:4}
Prove que o \omterm{def:b2:metric:def}{espaço métrico} discreto $X$ (conjunto qualquer) é \omterm{def:b2:metric:complete}{completo} e
que ele é \omterm{def:b2:metric:compact}{compacto} se e somente se $X$ é finito. Quais partes são
\omterm{def:b2:metric:connected}{conexas}?
\end{exercise}

\begin{exercise}[$\star\star$]\label{exo:b2:metric:5}
Sejam $X$ \omterm{def:b2:metric:compact}{compacto} e $f \colon X \to X$ com
\[
d\bigl(f(x), f(y)\bigr) < d(x, y) \quad \text{para todo } x \neq y .
\]
Prove que $f$ tem um único ponto fixo \emph{(minimize $x \mapsto
d(x, f(x))$)} e dê um exemplo em $X = \intco{1}{+\infty}$
(não \omterm{def:b2:metric:compact}{compacto}) sem ponto fixo.
\end{exercise}

\begin{exercise}[$\star\star$]\label{exo:b2:metric:6}
(\omterm{def:b2:metric:compact}{Compactos} encaixados) Seja $(K_n)$ uma sequência decrescente de partes
\omterm{def:b2:metric:compact}{compactas} não vazias de um \omterm{def:b2:metric:def}{espaço métrico}. Prove que $\bigcap_n K_n \neq
\emptyset$ \emph{(escolha $x_n \in K_n$ e extraia)}. Mostre por um exemplo
que fechados encaixados não vazios em $\R$ podem ter interseção
vazia.
\end{exercise}

\begin{exercise}[$\star\star$]\label{exo:b2:metric:7}
Sejam $K$ \omterm{def:b2:metric:compact}{compacto} e $f \colon K \to Y$ \omterm{def:b2:metric:continuity}{contínua} e bijetiva.
Prove que $f^{-1}$ é \omterm{def:b2:metric:continuity}{contínua} \emph{(use conjuntos fechados:
\cref{thm:b2:metric:globalcontinuity} e
\cref{thm:b2:metric:compactprops})}. Dê um contraexemplo sem
\omterm{def:b2:metric:compact}{compacidade} ($\gamma(t) = (\cos t, \sin t)$ em
$\intco{0}{2\pi}$).
\end{exercise}

\begin{exercise}[$\star\star$]\label{exo:b2:metric:8}
O \emph{conjunto de Cantor} $C$ obtém-se de $\intcc{0}{1}$
apagando repetidamente os terços médios \omterm{def:b2:metric:topology}{abertos}. Prove que $C$ é \omterm{def:b2:metric:compact}{compacto},
tem interior vazio e é infinito --- de fato \omterm{def:b2:structures:countable}{equipotente} a
$\{0,1\}^{\N}$ \emph{(expansões ternárias com algarismos $0,2$;
\cref{exo:b2:structures:3})}.
\end{exercise}

\begin{exercise}[$\star\star\star$]\label{exo:b2:metric:9}
Seja $X$ um \emph{\omterm{def:b2:metric:compact}{espaço métrico compacto}} e $f \colon X \to X$ uma
isometria: $d(f(x), f(y)) = d(x, y)$. Prove que $f$ é sobrejetiva.
\emph{Sugestão: se $a \notin f(X)$, então $\varepsilon = d(a, f(X)) > 0$
(por quê?); estude a órbita $a, f(a), f^2(a), \dots$ e mostre que seus
pontos estão dois a dois a distância $\geq \varepsilon$ --- contradição com a
\omterm{def:b2:metric:compact}{compacidade}.}
\end{exercise}

\begin{exercise}[$\star\star\star$]\label{exo:b2:metric:10}
Prove que $GL_n(\C)$ é \omterm{def:b2:metric:connected}{conexo} por caminhos. \emph{Sugestão: dada $A, B$
invertível, considere $p(z) = \det\bigl((1 - z)A + zB\bigr)$ para $z
\in \C$: um polinômio em $z$, não identicamente nulo, logo com
um número finito de raízes; ligue $0$ a $1$ em $\C$ por um caminho que as evite.}
\end{exercise}

\begin{exercise}[$\star\star$]\label{exo:b2:metric:11}
Para $\emptyset \neq A \subseteq X$, ponha $d(x, A) = \inf_{a \in A}
d(x, a)$. Prove que $x \mapsto d(x, A)$ é $1$-lipschitziana, que
$d(x, A) = 0$ se e somente se $x \in \overline A$, e que, para fechados não vazios
\emph{disjuntos} $A, B$, a função
\[
\varphi(x) = \frac{d(x, A)}{d(x, A) + d(x, B)}
\]
está bem definida, é \omterm{def:b2:metric:continuity}{contínua}, vale $0$ exatamente em $A$ e
$1$ exatamente em $B$ --- um ``interruptor'' \omterm{def:b2:metric:continuity}{contínuo} que separa dois
fechados disjuntos quaisquer.
\end{exercise}

\begin{exercise}[$\star\star\star$]\label{exo:b2:metric:12}
(Baire) Sejam $X$ um \omterm{def:b2:metric:complete}{espaço métrico completo} e $(U_n)_{n\geq1}$ uma
sequência de \omterm{def:b2:metric:topology}{abertos} densos. Prove que $\bigcap_n U_n$ é
denso em $X$ \emph{(dentro de uma bola qualquer, construa bolas fechadas encaixadas
$\overline B(x_n, r_n) \subseteq U_n$ com $r_n \to 0$ e use a
\omterm{def:b2:metric:complete}{completude})}. Deduza que $\R$ não é união \omterm{def:b2:structures:countable}{enumerável} de
fechados de interior vazio e recupere --- mais uma vez --- que
$\R$ não é \omterm{def:b2:structures:countable}{enumerável}.
\end{exercise}

\section{Problema: a iteração de Picard}

\omterm{def:b2:metric:complete}{Completude} mais contração é uma máquina de resolver: alimente-a com uma
equação escrita como problema de ponto fixo e ela devolve
existência, unicidade, um algoritmo e barras de erro. Este problema de fim
de semana roda a máquina a plena potência na equação $y' = f(t,
y)$: demonstramos o \emph{teorema local de Picard--Lindelöf} (o
coração não linear da teoria de Cauchy--Lipschitz do~\cref{ch:b2:diffeq}),
vemos cada hipótese ganhar o pão por meio de
contraexemplos e colhemos dividendos puramente métricos ---
a dependência \omterm{def:b2:metric:continuity}{contínua} dos dados, a equação de Kepler e a
autossemelhança do conjunto de Cantor.

\begin{omfigure}
\begin{tikzpicture}
\begin{axis}[omaxis, width=8.6cm, height=5.4cm,
    xmin=-0.05, xmax=1.45, ymin=0, ymax=4.2,
    xtick={0,1}, ytick={1,2,3}]
  \addplot[omDef!45, thick, domain=0:1.4, samples=2] {1};
  \addplot[omDef!65, thick, domain=0:1.4, samples=2] {1 + x};
  \addplot[omDef, thick, domain=0:1.4, samples=60]
    {1 + x + x^2/2};
  \addplot[omThm, very thick, domain=0:1.4, samples=120]
    {exp(x)};
  \node[font=\small, omDef] at (axis cs:1.28,0.8) {$y_0$};
  \node[font=\small, omDef] at (axis cs:1.28,2.15) {$y_1$};
  \node[font=\small, omDef] at (axis cs:1.28,2.95) {$y_2$};
  \node[font=\small, omThm] at (axis cs:1.17,3.85)
    {$\eu^{\,t}$};
\end{axis}
\end{tikzpicture}

{\small Iterados de Picard para $y' = y$, $y(0) = 1$: cada passagem
por $T(y)(t) = 1 + \int_0^t y$ acrescenta um termo de Taylor, e a
contração comprime toda a sequência uniformemente sobre
$\eu^{\,t}$.}
\end{omfigure}

\begin{problem}[{Problema de fim de semana --- o teorema de Picard--Lindelöf}]
\label{pb:b2:metric:1}
Em todo o problema, $t_0 \in \R$, $y_0 \in \R$, $a, b > 0$, e $f$ é uma
função \omterm{def:b2:metric:continuity}{contínua} no retângulo $R = \intcc{t_0 - a}{t_0 +
a} \times \intcc{y_0 - b}{y_0 + b}$, limitada por $M =
\sup_R\,\abs f$ e \emph{$L$-lipschitziana na segunda variável}:
$\abs{f(t, y) - f(t, z)} \leq L\abs{y - z}$ sempre que os dois pontos
estejam em $R$. Ponha
\[
h = \min\Bigl(a, \frac bM\Bigr) \quad (\text{com } h = a
\text{ se } M = 0), \qquad I = \intcc{t_0 - h}{t_0 + h}.
\]

\medskip
\noindent\textbf{Parte I --- O palco \omterm{def:b2:metric:complete}{completo}.}
\begin{enumerate}
  \item Demonstre os dois enunciados citados na
        \cref{def:b2:metric:complete}: uma parte fechada de um
        \omterm{def:b2:metric:complete}{espaço métrico completo} é \omterm{def:b2:metric:complete}{completa}, e uma parte \omterm{def:b2:metric:complete}{completa}
        de um \omterm{def:b2:metric:def}{espaço métrico} qualquer é fechada. Deduza que toda parte fechada
        de $\bigl(C(I), d_\infty\bigr)$ é um \omterm{def:b2:metric:complete}{espaço métrico
        completo}.
  \item Mostre que a aplicação $C(I) \to C(I)$, $y \mapsto \bigl(t
        \mapsto \int_{t_0}^{t}y(s)\,\dd s\bigr)$, é
        $h$-lipschitziana para $d_\infty$.
  \item (Os pontos fixos se movem menos que as aplicações) Seja $g \colon X
        \to X$ uma $k$-contração de um \omterm{def:b2:metric:def}{espaço métrico} com ponto
        fixo $\ell_g$, e seja $\widetilde g \colon X \to X$
        \emph{uma} aplicação qualquer com ponto fixo $\ell_{\widetilde g}$.
        Demonstre
        \[
        d(\ell_g, \ell_{\widetilde g}) \leq
        \frac{d\bigl(g(\ell_{\widetilde g}),
        \widetilde g(\ell_{\widetilde g})\bigr)}{1 - k}
        \leq \frac{\sup_{x \in X} d\bigl(g(x), \widetilde
        g(x)\bigr)}{1 - k}.
        \]
  \item (O truque do iterado) Seja $X$ \omterm{def:b2:metric:complete}{completo} não vazio e $g
        \colon X \to X$ uma aplicação --- não suposta \omterm{def:b2:metric:continuity}{contínua} ---
        tal que algum iterado $g^m$ seja uma $k$-contração.
        Prove que $g$ tem um único ponto fixo $\ell$ e que
        \emph{toda} órbita $x_{n+1} = g(x_n)$ converge para
        $\ell$. \emph{(Os pontos fixos de $g$ são pontos fixos de
        $g^m$; reciprocamente, $g(\ell)$ é ponto fixo de $g^m$;
        separe a órbita segundo os restos módulo $m$.)}
\end{enumerate}

\medskip
\noindent\textbf{Parte II --- O teorema de Picard--Lindelöf.}
\begin{enumerate}[resume]
  \item Mostre que uma função $y \colon I \to \intcc{y_0 - b}{y_0
        + b}$ é $C^1$ com $y(t_0) = y_0$ e $y' = f(t, y)$ em
        $I$ se e somente se ela é \omterm{def:b2:metric:continuity}{contínua} e satisfaz a
        equação integral
        \[
        y(t) = y_0 + \int_{t_0}^{t} f\bigl(s, y(s)\bigr)\dd s
        \qquad (t \in I).
        \]
  \item Sejam $X_h = \{y \in C(I) : \abs{y(t) - y_0} \leq b
        \text{ em } I\}$ e $T$ definida por $T(y)(t) =
        y_0 + \int_{t_0}^{t}f(s, y(s))\dd s$. Mostre que $X_h$ é
        uma parte fechada não vazia de $C(I)$, logo \omterm{def:b2:metric:complete}{completa}, e
        que $T$ leva $X_h$ em $X_h$ --- é aqui que $h \leq
        b/M$ trabalha.
  \item Mostre que $d_\infty\bigl(T(y), T(z)\bigr) \leq
        Lh\,d_\infty(y, z)$ em $X_h$: se $Lh < 1$, o teorema de
        Banach já conclui. Removemos essa condição de pequenez
        a seguir.
  \item Prove por indução sobre $n$:
        \[
        \abs{T^n(y)(t) - T^n(z)(t)} \leq
        \frac{\bigl(L\abs{t - t_0}\bigr)^n}{n!}\,
        d_\infty(y, z) \qquad (y, z \in X_h,\ t \in I),
        \]
        de modo que algum iterado de $T$ é uma contração. Conclua com a
        questão 4 (\emph{o teorema de Picard--Lindelöf}): o
        problema de Cauchy $y' = f(t,y)$, $y(t_0) = y_0$ tem exatamente
        uma solução em $I = \intcc{t_0 - h}{t_0 + h}$ com
        valores em $\intcc{y_0 - b}{y_0 + b}$.
  \item Mostre que a restrição ``com valores em $\intcc{y_0 -
        b}{y_0 + b}$'' é automática: toda solução do problema de
        Cauchy definida em $I$ cujo gráfico começa em $R$ permanece
        em $\intcc{y_0 - b}{y_0 + b}$ \emph{(considere o primeiro
        instante de saída e majore $\abs{y(t) - y_0}$ por $M\abs{t -
        t_0}$)}. Logo a unicidade vale entre todas as soluções em
        $I$.
  \item Rode a máquina em $y' = y$, $y(0) = 1$, partindo da
        constante $y^{(0)} \equiv 1$: calcule os iterados
        de Picard $y^{(n)}$, identifique-os e descreva a
        convergência.
\end{enumerate}

\medskip
\noindent\textbf{Parte III --- Cada hipótese ganha o pão.}
\begin{enumerate}[resume]
  \item (A condição \omterm{def:b2:metric:continuity}{lipschitziana} falha, a unicidade falha) Para $y' =
        2\sqrt{\abs y}$, $y(0) = 0$: verifique que $y \equiv 0$
        e que, para todo $c \geq 0$, a função $y_c(t) = 0$ para
        $t \leq c$, $y_c(t) = (t - c)^2$ para $t > c$, são todas
        soluções $C^1$ em $\R$. Onde exatamente $y \mapsto
        2\sqrt{\abs y}$ deixa de ser \omterm{def:b2:metric:continuity}{lipschitziana}?
  \item (A localidade é real) Para $y' = y^2$, $y(0) = 1$: resolva
        explicitamente, dê o intervalo maximal de existência e
        calcule o melhor $h$ que o teorema consegue certificar sobre todas as
        escolhas do retângulo ($a$ grande, $b$ livre): mostre que
        $h_{\max} = \sup_{b>0} \frac{b}{(1+b)^2} = \frac14$,
        enquanto a solução verdadeira vive em $\intoo{-\infty}{1}$.
  \item (A \omterm{def:b2:metric:complete}{completude} não é enfeite) Em $X = \Q \cap
        \intcc{1}{2}$ com a distância usual, seja $g(x) =
        \frac x2 + \frac1x$. Mostre que $g(X) \subseteq X$, que $g$
        é uma $\frac12$-contração \emph{(desigualdade do valor
        médio)} e que $g$ não tem ponto fixo em $X$.
        Qual hipótese do teorema de Banach falha, e qual é
        o ponto fixo no completado?
  \item (Barras de erro) Para uma $k$-contração $g$ num espaço
        \omterm{def:b2:metric:complete}{completo}, demonstre a estimativa \emph{a posteriori} $d(x_n,
        \ell) \leq \frac{k}{1-k}\,d(x_n, x_{n-1})$. Para a
        aplicação de Heron $g(x) = \frac x2 + \frac 1x$ em $\intcc{1}{2}$
        (ponto fixo $\sqrt2$), partindo de $x_0 = \frac32$:
        quantos passos a cota \emph{a priori}
        $\frac{k^n}{1-k}d(x_1, x_0)$ exige para uma precisão
        $10^{-6}$, e quantos passos bastam na realidade?
        (Calcule $x_1, x_2, x_3$ e seus erros; a cota da
        contração é honesta, mas pessimista ---
        Heron converge quadraticamente.)
\end{enumerate}

\medskip
\noindent\textbf{Parte IV --- Dependência \omterm{def:b2:metric:continuity}{contínua}.} Nesta
parte $Lh < 1$, de modo que o próprio $T$ é uma contração em $X_h$ (questão
7); escreva $y[\,y_0\,]$ para a solução de valor inicial
$y_0$.
\begin{enumerate}[resume]
  \item (Dependência do valor inicial) Seja $z_0$ outro
        valor inicial com $\abs{z_0 - y_0}$ pequeno o bastante para que
        os dois problemas caibam no retângulo. Usando a questão 3,
        demonstre
        \[
        d_\infty\bigl(y[y_0], y[z_0]\bigr) \leq
        \frac{\abs{y_0 - z_0}}{1 - Lh}.
        \]
  \item (Intervalos longos por encadeamento) Suponha que as soluções existam
        num segmento longo cortado em $m$ pedaços consecutivos em
        cada um dos quais a cota anterior se aplica com $Lh \leq
        \frac12$. Mostre que o desvio cresce por um fator no máximo
        $2$ por pedaço, logo $d_\infty \leq 2^m\abs{y_0 - z_0}$
        no total --- uma cota exponencial no comprimento, a sombra
        discreta do $\eu^{L\abs{t - t_0}}$ do lema de Gronwall
        (\cref{ch:b2:diffeq}).
  \item (Dependência do campo) Seja $g$ outro campo em
        $R$, também $L$-lipschitziano em $y$, com $\sup_R \abs{f - g}
        \leq \varepsilon$. Prove que as soluções
        correspondentes satisfazem $d_\infty \leq
        \frac{\varepsilon h}{1 - Lh}$: o erro de
        modelagem se propaga linearmente.
  \item (Parâmetros) Se uma família $f_\lambda$ de campos é
        $L$-lipschitziana em $y$ uniformemente e $\sup_R\abs{f_\lambda
        - f_\mu} \leq C\abs{\lambda - \mu}$, deduza que
        $\lambda \mapsto y_\lambda$ é \omterm{def:b2:metric:continuity}{lipschitziana} do
        espaço de parâmetros em $\bigl(C(I), d_\infty\bigr)$.
  \item (Sistemas não custam nada) Explique por que as Partes I, II e IV
        valem palavra por palavra para $y$ com valores em $\R^n$ (distâncias
        do sup construídas sobre qualquer uma das distâncias do
        \cref{ex:b2:metric:examples}), e depois calcule todos os iterados
        de Picard para o sistema $y' = Ay$, $y(0) = (c_1,
        c_2)$, $A = \begin{psmallmatrix} 0 & 1\\ 0 &
        0\end{psmallmatrix}$: mostre que a iteração se torna
        estacionária na solução exata após um passo.
\end{enumerate}

\medskip
\noindent\textbf{Parte V --- Dividendos métricos e síntese.}
\begin{enumerate}[resume]
  \item (Perturbação da identidade) Seja $X$ um palco \omterm{def:b2:metric:complete}{completo}
        do tipo espaço normado: tome $X = C(I)$ ou $\R^n$. Se
        $\eta \colon X \to X$ é $k$-lipschitziana com $k < 1$,
        prove que $x \mapsto x + \eta(x)$ é uma bijeção de $X$
        cuja inversa é $\frac1{1-k}$-lipschitziana \emph{(para cada
        $y$, aplique Banach a $x \mapsto y - \eta(x)$)}. Esse é
        o coração métrico do teorema da função inversa
        (\cref{ch:b2:diffcalc}).
  \item (Equação de Kepler) Para $0 \leq e < 1$ e $m \in \R$,
        prove que $x = m + e\sin x$ tem exatamente uma solução,
        que a iteração $x_{n+1} = m + e\sin x_n$ converge
        para ela a partir de qualquer início, e estime: para $e = \frac12$,
        $m = 1$, quantas iterações garantem um erro $\leq
        10^{-3}$ pela cota a priori? (A solução é $x
        \approx 1.4987$.)
  \item (O conjunto de Cantor é um ponto fixo) Sejam $S_1(x) = \frac
        x3$ e $S_2(x) = \frac x3 + \frac23$ em $\R$, e seja
        $C$ o conjunto de Cantor do~\cref{exo:b2:metric:8}. Prove que
        $C = S_1(C) \cup S_2(C)$ e explique em uma frase
        por que nenhum \emph{outro} \omterm{def:b2:metric:compact}{compacto} não vazio satisfaz essa
        equação (a aplicação $A \mapsto S_1(A) \cup S_2(A)$ é uma
        contração para uma distância entre \omterm{def:b2:metric:compact}{compactos} --- a
        distância de Hausdorff, tornada honesta no volume do terceiro ano de graduação).
  \item (A \omterm{def:b2:metric:connected}{conexidade} globaliza a unicidade) Seja $f$ localmente
        \omterm{def:b2:metric:continuity}{lipschitziana} em $y$ num \omterm{def:b2:metric:topology}{aberto}, e sejam $y, z$ duas
        soluções de $y' = f(t, y)$ num intervalo comum $J$
        com $y(t_0) = z(t_0)$. Prove que $y = z$ em $J$: mostre que
        $\{t \in J : y(t) = z(t)\}$ é não vazio, fechado em $J$
        e \omterm{def:b2:metric:topology}{aberto} em $J$ (pela unicidade local), e use a
        \omterm{def:b2:metric:connected}{conexidade} dos intervalos
        (\cref{thm:b2:metric:connectedness}).
  \item (Nada de pequenez para equações lineares) Para $y' =
        \alpha(t)y + \beta(t)$ com $\alpha, \beta$ \omterm{def:b2:metric:continuity}{contínuas}
        num segmento $\intcc{A}{B}$, adapte a questão 8 para mostrar que
        a cota fatorial vale no segmento \emph{inteiro},
        de modo que existência e unicidade são globais aí --- o
        caso escalar do teorema de Cauchy--Lipschitz do~\cref{ch:b2:diffeq},
        sem restrição alguma sobre o comprimento $B - A$.
  \item (Síntese) Uma frase para cada: o que a \omterm{def:b2:metric:complete}{completude}
        contribuiu; o que a contração contribuiu; o que o
        truque do iterado comprou em comparação com o Banach simples; onde entrou a
        \omterm{def:b2:metric:connected}{conexidade}; e qual contraexemplo da Parte
        III guarda qual hipótese. Nomeie o teorema-cume
        e diga o que substitui a contração quando $f$ é apenas
        \omterm{def:b2:metric:continuity}{contínua} (o teorema de Peano, via \omterm{def:b2:metric:compact}{compacidade} em espaços de
        funções --- o Arzelà--Ascoli do volume do terceiro ano de graduação).
\end{enumerate}
\end{problem}
